\section{Literature notation}\label{app:B.2}

The dark-matter literature and the condensed-matter literature use the same letters for different objects. The letter $\chi$ is the dark-matter particle in one and a susceptibility in the other, $q$ is the momentum transfer in one and the phonon wavevector in the other, and $W$ is the loss function, the screened interaction or the Debye--Waller exponent, depending on the paper. This book gives each letter one meaning, so several common symbols are renamed. \Cref{tab:B-literature} lists the renamings that a reader moving between the book and the papers will meet. The left column gives what the literature writes, the second column what this book writes, and the last column the section where the book's symbol is defined.

{\small
\setlength{\tabcolsep}{4pt}
\begin{longtable}{@{}>{\raggedright\arraybackslash}p{0.25\textwidth}>{\raggedright\arraybackslash}p{0.17\textwidth}>{\raggedright\arraybackslash}p{0.43\textwidth}>{\raggedright\arraybackslash}p{0.09\textwidth}@{}}
\caption{Symbols of the literature and of this book. Where a paper's symbol differs from the book's only by a factor or a sign, the factor is given in the Meaning column.}\label{tab:B-literature}\\
\toprule
Literature & Here & Meaning & Section \\
\midrule
\endfirsthead
\toprule
Literature & Here & Meaning & Section \\
\midrule
\endhead
\bottomrule
\endfoot
\multicolumn{4}{@{}l}{\itshape Response and screening}\\*
$\chi_{AB}$, $\chi$, $C^R_{AB}$ & $R_{ab}$ & retarded response function; the literature usually omits the factor $1/\Vol$ & \labelcref{sec:4.1} \\
$\Pi$ & $\tilde R$ & polarization, the Coulomb-proper (irreducible) response & \labelcref{sec:4.4} \\
$W$, $\Imag(-1/\varepsilon)$ & $L$ & loss function $\Imag[-(\varepsilon^{-1})_{00}]$ & \labelcref{sec:4.4} \\
$W$ (as in $GW$) & $W^{\mathrm{scr}}$ & screened Coulomb interaction & \labelcref{sec:4.8} \\
$-2\Imag G^{\mathrm{R}}$ & $\SpecA$ & spectral function; the book's $\SpecA$ integrates to one, the literature's to $2\pi$ & \labelcref{sec:4.7} \\
$G(\vect{r},t)$ & $\mathcal{G}_{\mathrm{VH}}(\vect{r},t)$ & Van Hove correlation function & \labelcref{sec:4.2} \\
$\lambda_{\mathrm{TF}}$ & $1/k_{\mathrm{TF}}$ & Thomas--Fermi screening length & \labelcref{sec:4.4} \\
$\kappa$ & $\varepsilon_{\mathrm{bg}}$ & background (remote-band) dielectric constant & \labelcref{sec:7.3} \\
$\alpha(\omega)$ & $\alpha_{\mathrm{abs}}(\omega)$ & optical absorption coefficient & \labelcref{sec:6.5} \\
\addlinespace
\multicolumn{4}{@{}l}{\itshape Dark matter and kinematics}\\*
$\varepsilon$, $\kappa$, $\chi$ & $\zeta$ & kinetic mixing of the dark photon & \labelcref{sec:1.4} \\
$g_\chi$, $g_e$ & $y_\chi$, $y_e$ & mediator couplings to dark matter and electrons & \labelcref{sec:1.4} \\
$n_\chi$ & $\nd_\chi$ & dark-matter number density $\rho_\chi/\mchi$ & \labelcref{sec:1.1} \\
$\vect{k}$ & $\vect{q}$ & momentum transferred from the dark matter to the target & \labelcref{sec:1.3} \\
$\dd\Omega$ & $\dd\Omega_{\vect{q}}$ & solid-angle element of the scattered momentum & \labelcref{sec:2.2} \\
$\sigma_e$ (fermionic absorption) & $\sigma_F$ & absorption cross-section of fermionic dark matter & \labelcref{sec:10.9} \\
$Q$ & $N_{\mathrm{eh}}$ & number of electron--hole pairs (ionization yield) & \labelcref{sec:1.5} \\
\addlinespace
\multicolumn{4}{@{}l}{\itshape Nuclei and phonons}\\*
$\vect{q}$, $\vect{k}$ & $\vect{Q}$ & phonon wavevector in the first zone & \labelcref{sec:5.2} \\
$r_n$, $R_1$; $s$ & $R_{\mathrm{nuc}}$; $s_{\mathrm{skin}}$ & Helm radius; skin thickness & \labelcref{sec:2.2} \\
$c$, $a$ & $c_{\mathrm{nuc}}$, $a_{\mathrm{nuc}}$ & Woods--Saxon half-density radius and diffuseness & \labelcref{sec:2.2} \\
$\lambda$, $\mu$ (Lam\'e) & $C_{12}$, $C_{44}$ & elastic constants of an isotropic medium & \labelcref{app:C.18} \\
$\lambda$ & $\lambda_{\mathrm{ep}}$ & electron--phonon coupling constant & \labelcref{sec:6.4} \\
$\mu^*$ & $\mu^*_{\mathrm{C}}$ & Coulomb pseudopotential & \labelcref{app:E.9} \\
$Z$, $\Delta$ & $Z^{\mathrm{El}}$, $\Delta^{\mathrm{El}}$ & Eliashberg renormalization function and gap & \labelcref{app:E.9} \\
\addlinespace
\multicolumn{4}{@{}l}{\itshape Electrons in crystals}\\*
$\mu$ & $\mu_{\mathrm{ch}}$ & chemical potential & \labelcref{sec:2.5} \\
$N(0)$, $N(E_{\mathrm{F}})$ & $\mathcal{N}_{\mathrm{F}}$ & density of states per spin and volume at the Fermi level & \labelcref{sec:2.5} \\
$v_{\mathrm{F}}$ (Dirac material) & $v_{\mathrm{D}}$ & Dirac velocity; $v_{\mathrm{F}}$ is the Fermi velocity of the electron gas & \labelcref{sec:7.3} \\
$\chi_{nl}$ & $\tilde\psi_{nl}$ & radial momentum-space wavefunction of an atomic shell & \labelcref{sec:2.2} \\
$T_{\mathrm{s}}$ & $K_{\mathrm{s}}$ & Kohn--Sham (independent-electron) kinetic energy & \labelcref{sec:3.4} \\
$V_2$ & $V_{\mathrm{cov}}$ & covalent energy between two bond hybrids & \labelcref{sec:3.2} \\
$\Omega$; $\Omega_{\mathrm{I}}$; $\Omega_{\mathrm{D}} + \Omega_{\mathrm{OD}}$ & $\Omega_W$; $\Omega_W^{\mathrm{I}}$; $\tilde\Omega_W$ & Wannier spread and its gauge-invariant and gauge-dependent parts & \labelcref{sec:3.6} \\
$V_{\mathrm{cell}}$ & $\Omega$ & unit-cell volume & \labelcref{sec:3.1} \\
$\chi$, $\chi^0$, $\alpha$ & $R^U$, $R^{U,0}$, $\delta V^U$ & interacting and bare occupation responses, and the potential shift, of linear-response Hubbard $U$ & \labelcref{app:E.9} \\
$\alpha$, $K$ & $\alpha_{\mathrm{ohm}}$ & ohmic coupling of a tunneling defect to conduction electrons & \labelcref{sec:16.8} \\
$e$ (Gaussian) & $e_{\mathrm{G}}$ & Gaussian elementary charge, $e_{\mathrm{G}}^2 = \alpha$; the book's $e$ has $e^2 = 4\pi\alpha$ & \labelcref{app:A.2} \\
$k$ (elliptic modulus) & $s$ & modulus of the complete elliptic integrals & \labelcref{app:C.8} \\
\end{longtable}
}

The Fourier sign, the sign of the mixing term, the units of charge and the normalization of the structure factor differ from parts of the literature by a sign or a factor rather than by a letter, and a formula taken from a paper must be converted before it is combined with the book's.

\emph{Fourier transforms.} A function of position is transformed with $\ee^{+\ii\vect{q}\cdot\vect{r}}$ and a function of time with $\ee^{+\ii\omega t}$, with the factor $(2\pi)^{-1}$ per dimension on the inverse transform (\cref{eq:A.3-space,eq:A.3-time}). The electron density operator is then $n_e(\vect{q}) = \sum_i\ee^{+\ii\vect{q}\cdot\vect{r}_i}$, with $\vect{q}$ the momentum given to the target. Some papers use the opposite spatial sign with the same direction of $\vect{q}$ and write the density operator as $\sum_i\ee^{-\ii\vect{q}\cdot\vect{r}_i}$~\cite[Eq.~(S.4)]{hochberg2021determining}; others keep $\ee^{+\ii\vect{q}\cdot\vect{x}}$ on the target by evaluating the transform of the potential at $-\vect{q}$~\cite[Eq.~(4)]{trickle2022effective}. A perturbation of frequency $\omega > 0$ goes as $\ee^{-\ii\omega t}$, so retarded functions are analytic for $\Imag\omega > 0$ and $\varepsilon_2 > 0$ in an absorbing medium. Engineering texts use $\ee^{+\ii\omega t}$, and their dielectric functions are the complex conjugates of ours (\cref{app:A.3}).

\emph{Sign of the mixing term.} The dark photon mixes through $+(\zeta/2)F_{\mu\nu}F'^{\mu\nu}$, with $\mathcal{L} \supset -eA_\mu J^\mu_{\mathrm{em}}$ for the photon, so that the dark photon couples to electrons with $y_e = -\zeta e$ (\cref{sec:1.4,app:A.9}). Ref.~\cite[Eq.~(70)]{lin2019tasi} has the same mixing sign but writes $\kappa$ for $\zeta$ and $+eA_\mu J^\mu_{\mathrm{em}}$; Ref.~\cite[Eq.~(2.1)]{essig2016direct} writes $\varepsilon$ and mixes with hypercharge, $(\varepsilon/2\cos\theta_W)F_{Y\mu\nu}F'^{\mu\nu}$, with $\theta_W$ the weak mixing angle. Changing either sign flips the signs of $y_e$ and $\tilde V$. No rate changes, because every rate contains $\zeta^2$.

\emph{Units of charge.} Charges are in Heaviside--Lorentz units, $e^2 = 4\pi\alpha$, so the Coulomb energy of two electrons is $e^2/(4\pi r)$ and the Coulomb kernel is $v_{\mathrm{C}}(k) = e^2/k^2$ (\cref{app:A.2}). Condensed-matter sources usually use Gaussian units, with $e_{\mathrm{G}}^2 = \alpha$ and kernel $4\pi e_{\mathrm{G}}^2/k^2$. A formula is converted by replacing every $e^2$ of this book with $4\pi e_{\mathrm{G}}^2$ for Gaussian units, or with $e_{\mathrm{SI}}^2/\varepsilon_0^{\mathrm{SI}}$ for SI units. The dimensionless response functions $\varepsilon$, $\varepsilon^{-1}$ and $L$ have the same values in all three systems.

\emph{Normalization of the structure factor.} The dynamic structure factor carries $2\pi/\Vol$ (\cref{eq:A.7-S}). It is intensive, has dimension $\mathrm{eV}^2$ for the density, and obeys the f-sum rule $\int\dd\omega\,\omega\,S(\vect{q},\omega) = \pi\nd_eq^2/m_e$ over the whole real line (\cref{eq:A.7-fsum}). Many-body texts define the structure factor per electron and without the $2\pi$; at $T = 0$ theirs is $S/(2\pi\nd_e)$, with first moment $q^2/2m_e$ (\cref{app:A.7}). Ref.~\cite[Eqs.~(2) and (4)]{knapen2021dark} uses the book's normalization, with $\chi$ for $R$ and $\vect{k}$ for the momentum transfer. For a spin-degenerate band structure the book's $S$ is summed over spin: its independent-particle form carries $4\pi/\Vol = 2\times2\pi/\Vol$ outside a sum over spatial bands (\cref{eq:A.7-spin}). Refs.~\cite[Eq.~(A.29)]{essig2016direct} and~\cite[Eq.~(6)]{catena2021crystal} also write the factor 2 outside their band sums, and Ref.~\cite[Appendix~B]{knapen2021dark} leaves it implicit in its main-text formulas. What differs between sources is whether the quantity called $S$, or the crystal form factor, includes the 2, and a factor of 2 between two published rates for the same material often comes from this.
